\documentclass[10pt,twocolumn]{article}
\usepackage[T1]{fontenc}
\usepackage[utf8]{inputenc}
\usepackage{lmodern}
\usepackage[margin=1.8cm,top=2.4cm,bottom=2cm,columnsep=0.6cm]{geometry}
\usepackage{fancyhdr}
\usepackage{titlesec}
\usepackage{booktabs}
\usepackage{amsmath,amssymb,amsfonts}
\usepackage{microtype}
\usepackage{xcolor}
\usepackage{mdframed}
\usepackage{parskip}
\usepackage{enumitem}
\usepackage{hyperref}

\definecolor{rulered}{RGB}{180,30,30}
\definecolor{boxgray}{RGB}{245,245,245}
\definecolor{keywordgray}{RGB}{100,100,100}

\pagestyle{fancy}
\fancyhf{}
\fancyhead[L]{\textsc{\small Claude Mag}}
\fancyhead[C]{\small\textit{Machine Learning Research Digest}}
\fancyhead[R]{\small 2026-08-05}
\fancyfoot[C]{\small\thepage}
\renewcommand{\headrulewidth}{0.8pt}

\titleformat{\section}{\normalsize\bfseries\color{rulered}}{}{0pt}{}%
  [\vspace{-4pt}\color{rulered}\rule{\columnwidth}{0.4pt}]
\titlespacing{\section}{0pt}{8pt}{4pt}

\hypersetup{colorlinks=true,urlcolor=rulered,linkcolor=black}

\begin{document}
\setlength{\parindent}{0pt}
\setlength{\parskip}{4pt}

{\small\color{keywordgray}\textbf{Keywords:}
  rectified flow \textbullet{}
  optimal transport \textbullet{}
  flow matching \textbullet{}
  statistical learning theory}\par\vspace{4pt}

{\LARGE\bfseries\raggedright
  Computational and Statistical Guarantees of the $c$-Rectified Flow}\par\vspace{4pt}

{\large\itshape\raggedright
  A Cost-Aware Flow Matching Variant Provably Minimises User-Specified Transport
  Costs While Preserving Marginals, With Sample Complexity Bounds}\par\vspace{6pt}

{\small\textbf{By} Leda Wang, Zhehao Xu, Qiang Liu \& Harrison H. Zhou
  (UT Austin; Yale University)}\par\vspace{2pt}

{\small\color{keywordgray}arXiv:2608.02487 \textbullet{} August 2026}
\par\vspace{2pt}\noindent{\color{rulered}\rule{\columnwidth}{0.6pt}}\vspace{6pt}

Flow matching trains fast generative models but offers no guarantee that the
learned transport minimises any user-specified cost. The \textit{$c$-Rectified Flow}
restricts learned vector fields to gradient-of-convex-conjugate form and proves
that iterative application monotonically reduces $c$-transport cost while
automatically preserving marginal distributions---with sample complexity bounds
characterising how many data points are needed.

\begin{mdframed}[backgroundcolor=boxgray,linecolor=rulered,linewidth=1pt,
    innerleftmargin=6pt,innerrightmargin=6pt,
    innertopmargin=6pt,innerbottommargin=6pt]
\textbf{\small AT A GLANCE}\par\vspace{4pt}
\begin{description}[leftmargin=0pt,labelsep=4pt,itemsep=2pt,parsep=0pt,
    font=\small\bfseries]
  \item[Problem:] Standard rectified flow / flow matching learns a valid generative
    model but provides no guarantee of minimising any user-specified transport cost.
  \item[Why it matters:] Many scientific applications (molecule generation,
    image editing, optimal control) have explicit cost functions; a theoretically
    grounded cost-aware flow model is essential for safe deployment.
  \item[Method:] Restrict vector fields to $\nabla\bar{c}(\nabla f_t(\cdot))$ form;
    iterative application provably reduces $c$-transport cost.
  \item[Result:] Computational and statistical guarantees for cost-aware flow learning;
    conditions for recovery of the true optimal transport solution.
  \item[Evidence:] Formal proofs of convergence and sample complexity; counterexamples
    delineating when strong assumptions are necessary.
\end{description}
\end{mdframed}

\section{The Big Picture}

Rectified flow and its generalisations (flow matching, stochastic interpolants)
have become the standard training framework for generative models used in image
synthesis, video generation, and scientific modelling. The core idea is simple:
train a neural ODE velocity field $v_t(x)$ to push samples from a source distribution
$\pi_0$ (e.g., Gaussian noise) towards a target distribution $\pi_1$ (e.g., data),
with training implemented as a simple regression objective.

This works remarkably well in practice. But for applications where the \emph{how}
of transport matters---molecule generation (minimise conformational energy),
image editing (preserve perceptual structure), optimal control (minimise actuator
effort)---standard flow matching provides no guarantee that the learned transport
is cost-efficient. It may find a valid coupling $\pi_0 \to \pi_1$ that is wasteful
or even physically unrealisable.

$c$-Rectified Flow closes this gap.

\section{Why Existing Approaches Fall Short}

\textbf{Standard rectified flow} minimises the squared-$L_2$ distance between a
straight-line trajectory and the learned trajectory, which corresponds to the
$L_2$ transport cost. It provides no guarantees for other cost functions and even
for $L_2$ it recovers the OT solution only under strong regularity conditions.

\textbf{Iterative rectification} (reflow) applies rectified flow recursively on
coupled samples to straighten trajectories, reducing $L_2$ cost in practice but
providing no monotone convergence guarantee for general costs $c$.

\textbf{Minibatch OT} pre-computes approximate couplings before training, but the
approximation error accumulates and there is no learning-theoretic analysis of
the resulting flow.

\section{Core Mechanism}

For a convex cost function $c: \mathbb{R}^d \times \mathbb{R}^d \to \mathbb{R}$,
let $\bar{c}$ denote its convex conjugate (Legendre--Fenchel transform):
\[
\bar{c}(u) = \sup_{x}\, \langle u, x \rangle - c(x).
\]

$c$-Rectified Flow restricts the learned velocity field to the structured form:
\[
v_t(x) = \nabla\bar{c}\!\left(\nabla f_t(x)\right),
\]
where $f_t: \mathbb{R}^d \to \mathbb{R}$ is an unconstrained time-dependent
scalar potential. The key observation is that any vector field of this form
is automatically the subdifferential of a convex function (since $\bar{c}$
is convex and $f_t$ is smooth), which is precisely the structure of the
optimal transport map for cost $c$ by Brenier's generalised theorem.

The neural network learns $f_t$ (unconstrained), and $v_t$ is computed from it
via automatic differentiation. The constraint is enforced implicitly by the
parameterisation, not by projection or regularisation.

\section{Technical Formulation}

\textbf{Training objective.}
Given source samples $\{x_0^{(i)}\}_{i=1}^n \sim \pi_0$ and target samples
$\{x_1^{(j)}\}_{j=1}^n \sim \pi_1$, paired by a coupling $\gamma_t$ at
interpolation time $t \in [0,1]$, the $c$-Rectified Flow objective is:
\[
\mathcal{L}_{c\text{-RF}}(\theta) =
\mathbb{E}_{(x_0, x_1) \sim \gamma}\!\left[
\left\| v_\theta(x_t, t) - \dot x_t \right\|^2
\right],
\quad x_t = (1-t)x_0 + t x_1,
\]
where $v_\theta = \nabla\bar{c}(\nabla f_\theta)$ is enforced by parameterisation.

\textbf{Computational guarantee.}
Let $\mathcal{T}_c(\gamma)$ denote the $c$-transport cost of coupling $\gamma$.
Theorem~1 in the paper establishes:
\[
\mathcal{T}_c(\gamma^{(k+1)}) \leq \mathcal{T}_c(\gamma^{(k)})
\]
for all $k$, where $\gamma^{(k)}$ is the coupling induced by the $k$-th iterate
of $c$-rectification. Marginals are preserved exactly at each iteration.

\textbf{Statistical guarantee.}
Theorem~2 provides: with $n$ i.i.d.\ samples from $\pi_0 \times \pi_1$, the
estimated $c$-transport cost satisfies:
\[
\mathcal{T}_c(\hat\gamma_n) - \mathcal{T}_c(\gamma^*) \leq
C_s \cdot n^{-2/(d+4)} \cdot \log n
\]
with high probability, where $C_s$ depends on the smoothness class of the
optimal transport map and the Lipschitz constant of $\bar{c}$.

\section{Learning or Inference Procedure}

\textbf{Iterative $c$-rectification:}
\begin{enumerate}[itemsep=2pt,parsep=0pt]
  \item Initialise coupling $\gamma^{(0)} = \pi_0 \otimes \pi_1$ (independent).
  \item For $k = 0, 1, 2, \ldots$:
    \begin{enumerate}[itemsep=1pt,parsep=0pt]
      \item Pair samples from $\pi_0$ and $\pi_1$ by simulating the current
        flow $v^{(k)}$ to obtain coupled pairs $(x_0, x_1^{(k)})$.
      \item Train $v^{(k+1)}$ by minimising $\mathcal{L}_{c\text{-RF}}$ on
        these coupled pairs.
    \end{enumerate}
\end{enumerate}

\textbf{Inference:} Simulate the learned ODE
$\dot{x}_t = v^{(K)}_\theta(x_t, t)$ from $x_0 \sim \pi_0$ to $x_1$
using a numerical ODE solver (e.g., Euler with $T$ steps).

\section{What the Guarantee Says}

The computational guarantee ensures that repeated $c$-rectification is a
monotone interior algorithm: each iterate reduces the cost without leaving
the feasible set of couplings with the correct marginals. This is analogous
to coordinate descent on a convex objective but here the geometry is the
space of probability couplings.

The statistical guarantee implies that at dimension $d$ and smoothness $s$,
the estimation error decays as $n^{-s/(2s+d)}$ (classical non-parametric rate),
showing that $c$-Rectified Flow is statistically efficient.

However, recovery of the true OT solution requires additional assumptions:
connected support of $\pi_0, \pi_1$ and smoothness of the OT map $T^*$.
The paper provides counterexamples showing these conditions are not vacuous.

\section{Experimental Findings}

\begin{itemize}[itemsep=2pt,parsep=0pt]
  \item On synthetic 2-D optimal transport problems with non-$L_2$ costs
    (e.g., $c(x,y) = \|x-y\|_1$, earth-mover cost), $c$-Rectified Flow
    achieves provably lower transport cost than standard flow matching after
    the same number of training iterations.
  \item Monotone cost decrease is confirmed empirically for $K = 1, 2, 3$
    iterations of $c$-rectification across five experimental settings.
  \item Sample complexity curves match the theoretical $n^{-2/(d+4)}$ rate
    on synthetic data with known ground-truth OT maps.
  \item On a molecular conformer generation task with steric clash cost,
    $c$-Rectified Flow produces conformers with lower energy than standard
    flow matching baselines.
\end{itemize}

\section{Ablations and Interpretation}

\textbf{Effect of $\bar{c}$ parameterisation:} Using a learned neural $\bar{c}$
rather than the closed-form convex conjugate produces slightly better results
but is harder to optimise and risks losing the theoretical guarantees if the
learned function is not truly convex.

\textbf{Iteration count $K$:} Cost decreases most sharply in the first iteration
($k = 0 \to 1$) and yields diminishing returns beyond $K = 3$. This matches the
theoretical prediction that later iterates are already close to the OT solution.

\textbf{Disconnect: theory vs.\ practice:} In settings without smooth OT maps
(e.g., multimodal targets with disconnected support), $c$-Rectified Flow converges
to a local cost minimum rather than the global OT solution, consistent with the
counterexamples in the theoretical analysis.

\section*{Reference}

Leda Wang, Zhehao Xu, Qiang Liu, Harrison H. Zhou.
\textbf{Computational and Statistical Guarantees of the $c$-Rectified Flow}.
arXiv:2608.02487, August 2026.
\url{https://arxiv.org/abs/2608.02487}

\end{document}
